\documentclass[a4paper,12pt]{jsarticle}
\usepackage[margin=25mm]{geometry}
\usepackage{amsmath,amssymb}
\usepackage{enumitem}

\begin{document}

\begin{center}
{\LARGE 数I・A 共通テスト本番レベル（大問4題）}
\end{center}

\vspace{1em}

\section*{問題}

\section*{第1問（数と式・二次関数）}
実数 \(a\) を定数とし、
\[
f(x)=x^2-2(a+1)x+a^2-2a+5
\]
とする。
\begin{enumerate}[label=\arabic*.]
  \item \(f(x)\) の最小値を \(a\) を用いて表せ。
  \item 不等式 \(f(x)\le 0\) が実数解をもつような \(a\) の範囲を求めよ。
  \item 方程式 \(f(x)=0\) が異なる2つの正の実数解をもつような \(a\) の範囲を求めよ。
  \item 3.のとき、2解を \(\alpha,\beta\)（\(\alpha<\beta\)）とするとき、\(\alpha<2<\beta\) となる \(a\) の範囲を求めよ。
\end{enumerate}

\section*{第2問（場合の数・確率）}
1から10までの整数が1枚ずつ書かれた10枚のカードから、同時に4枚を取り出す。
\begin{enumerate}[label=\arabic*.]
  \item 全事象は何通りか。
  \item 4枚の和が偶数となる確率を求めよ。
  \item 4枚のうち、3の倍数がちょうど1枚である確率を求めよ。
  \item 4枚の最大値と最小値の差が5である確率を求めよ。
\end{enumerate}

\section*{第3問（図形と計量）}
三角形 \(ABC\) において、
\[
AB=8,\quad AC=10,\quad \angle A=60^\circ
\]
とする。
\begin{enumerate}[label=\arabic*.]
  \item 辺 \(BC\) の長さを求めよ。
  \item 三角形 \(ABC\) の面積を求めよ。
  \item 三角形 \(ABC\) の外接円の半径 \(R\) を求めよ。
  \item 辺 \(BC\) を底辺としたときの高さを求めよ。
\end{enumerate}

\section*{第4問（データの分析）}
ある40人学級の数学テスト（100点満点）について、階級別度数が次の通りであった。
\[
\begin{array}{c|ccccc}
\text{階級(点)} & 0\!\sim\!20 & 20\!\sim\!40 & 40\!\sim\!60 & 60\!\sim\!80 & 80\!\sim\!100\\ \hline
\text{度数(人)} & 2 & 6 & 10 & 14 & 8
\end{array}
\]
また、このテストの標準偏差は \(18\) 点とする。
\begin{enumerate}[label=\arabic*.]
  \item 60点以上の相対度数を求めよ。
  \item 各階級の階級値を用いて平均点を近似し、その値を求めよ。
  \item 偏差値 \(H\) を
  \[
  H=50+10\frac{x-\bar{x}}{18}
  \]
  （\(\bar{x}\)：2.で求めた平均値）で定める。\(x=78\) のときの偏差値を求めよ。
  \item 偏差値が65以上となるための得点 \(x\) の範囲を求めよ。
\end{enumerate}

\end{document}
